\begin{abstract}
\noindent
An encrypted mempool lets block builders order transactions before they can read them. The transactions must then be opened by some party: a committee adds trust assumptions for both privacy and liveness, and a sender that opens its own transaction can withhold its key when the ordering is unfavorable. We present \leanseal, in which the sender opens its own transaction but attaches a time-lock puzzle built from SHA-256 alone, so that anyone can open the transaction without the sender after a fixed amount of sequential computation, chosen to exceed the time needed to fix the order: about 30 seconds on Ethereum. For a well-formed puzzle, withholding the key then cannot keep the transaction hidden, without a committee or a trusted setup; it can still delay or cancel the transaction's execution, which a deposit must deter.

The puzzle is the chained-hash puzzle of Mahmoody, Moran and Vadhan (CRYPTO 2011). For our encoding we prove, in the random-oracle model, that a classical attacker with the hash rate of the whole Bitcoin network derives the key early with probability at most about $2^{-117}$ for a 30-second lock. Against an attacker who computes one chained SHA-256 every 3.5\,ns, a 30-second puzzle costs the sender about 21\,s on a laptop CPU, can be built before the transaction exists, and adds 152 to 512 bytes to the transaction; opening with the revealed key takes 12\,ms. We also identify what \leanseal cannot provide: like every time lock, it also discloses transactions that are never included. We show how deposits, validity windows and composition with a committee mitigate its limitations, and compare it with threshold, TEE and hybrid designs.

\medskip
\noindent\textbf{Code, specification and benchmark logs:} \repo
\end{abstract}
